%%  Chapter 5 - Discussion

\ifpdf
\graphicspath{{chapters/Figs/}{chapters/Figs/Raster/}{Figs/PDF/}{Figs/}}
\else
\graphicspath{{chapters/Figs/}{chapters/Figs/Vector/}{Figs/}}
\fi

\chapter{Discussion}
\label{ch:discussion}

\section{The BOE Mechanism in Indian Monsoon Depressions}
\label{sec:disc_boe}

The convergent lines of evidence assembled in
Chapter~\ref{ch:results} support a coherent mechanistic narrative.
Long-lived MDs maintain inner-core LHF at levels comparable to the
Bay of Bengal after landfall (Section~\ref{subsec:hovmoller_slhtf}),
their boundary layers remain near-saturated, and soil moisture
increases progressively beneath the inner core
(Section~\ref{subsec:hovmoller_swvl1}).

\subsection{LHF Magnitude and Cloud Cover Limitations}
The LHF magnitudes observed (70--120\,W\,m$^{-2}$) are lower than
extreme Atlantic BOE cases ($>200$\,W\,m$^{-2}$) reported by
\autocite{andersen2014insights}. This quantitative difference is
physically expected. In the Indian summer monsoon, persistent cloud
cover reduces solar insolation at the surface compared to the clearer
skies of the US Great Plains. However, because the Indian soil is
often near-saturated (ASMI $> 0.35$), the Bowen ratio remains
exceptionally low. This suggests that in the Indian context, the BOE
is driven by the \emph{efficiency} of moisture transfer rather than
high absolute energy input.

% \subsection{Soil Moisture Memory and ASMI Thresholds}
% Our finding that a 7-day ASMI lag better predicts persistence than
% a 14-day lag aligns with the soil moisture memory timescales of
% clay-loam soils common in the Odisha and Chhattisgarh basins. The
% threshold of 0.35\,m$^3$m$^{-3}$ identified in
% Section~\ref{subsec:asmi_results} serves as a "hydrological
% trigger" for the BOE, echoing the proxy findings of \autocite{kishtawal2013}.

\section{The Eastern Ghats as a Thermodynamic Filter}
\label{sec:disc_ghats}

As shown in Section~\ref{sec:results_elevation}, the Eastern Ghats do not
act as a binary wall. Instead, they function as a filter. Non-BOE
systems undergo structural collapse due to orographic moisture
stripping on the windward slopes. Conversely, BOE-active systems
exploit the high leeward soil moisture. The leeward speed recovery
observed in the composite analysis (Section~\ref{subsec:ts_speed}) is
a significant operational finding: it suggests that the land surface
can re-energize a system after terrain-induced disruption, provided
the land-atmosphere coupling is sufficiently strong.

% \chapter{Discussion}
% \label{ch:discussion}
%
% \section{The BOE Mechanism in Indian Monsoon Depressions}
% \label{sec:disc_boe}
%
% The convergent lines of evidence assembled in
% Chapter~\ref{ch:results} support a coherent mechanistic narrative.
% Long-lived MDs maintain inner-core LHF at oceanic levels after
% landfall (Section~\ref{subsec:hovmoller_lhf}), their boundary layers
% remain near-saturated throughout the land phase
% (Section~\ref{subsec:hovmoller_dpd}), soil moisture increases
% progressively beneath the inner core
% (Section~\ref{subsec:hovmoller_swvl1}), and the Bowen ratio stays low
% (Section~\ref{subsec:ts_bowen}), all consistent with the BOE
% mechanism. These signals are absent or rapidly collapsing in
% short-lived systems.
%
% The LHF magnitudes observed in the Indian MD context (70 -
% 120\,W\,m$^{-2}$ during the land phase) are lower than those reported
% for extreme Atlantic BOE cases ($>200$\,W\,m$^{-2}$;
% \citealt{andersen2014insights, wakefield2021}). This quantitative
% difference is physically expected. Atlantic BOE events typically
% occur in late summer (August - October) over the warm, irrigated, or
% recently flooded soils of the US South and Midwest, where soil
% surface temperatures can reach 30 - 35$^\circ$C. Indian MDs traverse
% soils that are cooler (owing to monsoon cloud shading reducing
% insolation) and where the available radiation for evapotranspiration
% is partly offset by the thick cloud layer. Nevertheless, the
% 70\,W\,m$^{-2}$ BOE threshold derived by \autocite{andersen2013} is
% shown here to be achievable and sustained in the Indian context,
% confirming that the BOE is not an exclusively Atlantic phenomenon but
% operates in the Indian monsoon environment under appropriate soil
% moisture conditions.
%
% \subsection{The Self-Wetting Feedback and Its Upper Bound}
% \label{subsec:disc_feedback}
%
% The co-evolution of \texttt{swvl1} and precipitation in
% Section~\ref{subsec:ts_sm_tp} reveals that the BOE is not simply a
% response to antecedent conditions but involves an active
% self-reinforcing process during the land phase. This has an important
% implication: even a system making landfall with moderately low ASMI
% ($\sim 0.28$ - $0.35$\,m$^3$\,m$^{-3}$) may approach BOE-sustaining
% soil moisture levels after 12-18\,h of its own heavy precipitation.
% The antecedent condition determines the \emph{initial} LHF level at
% landfall; whether a system can sustain LHF above the BOE threshold
% over the subsequent 18 hours is also determined by how rapidly the
% self-wetting feedback can drive \texttt{swvl1} to the energy-limited regime.
%
% This creates a practical implication for forecasting: a system making
% landfall over moderately moist soil may still qualify as a BOE
% candidate if its own initial precipitation is heavy enough to
% saturate the surface quickly. Conversely, an initially dry surface
% may prevent the self-wetting from occurring fast enough, even if
% large-scale atmospheric moisture advection is available.
%
% \subsection{Comparison with Prior Indian MD Studies}
% \label{subsec:disc_prior_india}
%
% The present study directly extends and observationally validates a
% sequence of prior works.
% \autocite{kishtawal2013enhancement} demonstrated, using proxy
% precipitation as an antecedent moisture indicator, that BOB
% depression inland penetration correlates with pre-storm soil wetness.
% This study replaces the proxy with direct ERA5 \texttt{swvl1}
% diagnostics and confirms the relationship over a 24-year climatology.
% The mechanistic link which \autocite{kishtawal2013enhancement} could not
% directly demonstrate from their proxy-based approach is established
% here through the Hovm\"{o}ller LHF signal and the Bowen ratio time series.
%
% \autocite{hunt2017effect} and \autocite{osuri2020} demonstrated through NWP
% sensitivity experiments that soil moisture initialisation affects MD
% simulation. The present observational evidence complements these
% modelling results by showing that the ERA5 soil moisture state
% without NWP forward integration is itself predictive of MD longevity,
% confirming that the signal is robustly present in the reanalysis record.
%
% \autocite{chang2009} identified land surface feedback as a modulator of
% post-landfall MD precipitation structure in WRF simulations. The
% Hovm\"{o}ller precipitation patterns in
% Section~\ref{subsec:hovmoller_tp} provide the observational
% equivalent: the sustained inner-core and contracted outer rainband
% structure after landfall in BOE-candidate systems mirrors precisely
% the model-simulated precipitation
% structure that \autocite{chang2009} attributed to land surface feedback.
%
% \autocite{ankur2022} showed that irrigated land cover enhances soil
% moisture and modulates MD behaviour. This study implicitly captures
% this effect through the ASMI (which integrates all moisture inputs,
% including irrigation) but cannot separate irrigation-driven soil
% moisture from monsoon-rainfall-driven soil moisture. This remains a
% limitation and an opportunity for future work
% (Section~\ref{sec:disc_limitations}).
%
% The BOE framework applied by \autocite{vijay2024boe} to the 1999 Odisha
% Super Cyclone demonstrated BOE-like inland maintenance of a
% post-landfall system in the Indian Ocean basin. The present
% climatological analysis shows that this mechanism is not confined to
% extreme events like the 1999 Super Cyclone but is a recurring process
% across the annual MD population.
%
% \section{The Role of the Eastern Ghats and Its Interaction with the BOE}
% \label{sec:disc_ghats}
%
% The differential translational speed recovery documented in
% Section~\ref{subsec:ts_speed}
% provides the most direct evidence for a BOE-terrain interaction. Both
% long-lived and short-lived systems decelerate over the Ghats, but
% only long-lived (BOE-candidate) systems recover speed on the leeward
% side. This differential is interpreted as follows.
%
% The BOE maintains the pressure gradient force of the MD vortex by
% sustaining deep convection on the leeward side through continuous LHF
% from the wet soil.
% In the absence of BOE, frictional decoupling during Ghats crossing
% increases vortex tilt and degrades the efficiency of the convective
% organisation.
% Without a local surface energy source to sustain convection on the
% leeward side, the vortex cannot recover its coherence and the system
% dissipates. With BOE, the leeward soil (which has been wetted by the
% windward-side precipitation that spilled over the terrain gaps)
% provides sufficient LHF to re-moisten the boundary layer, maintain
% CAPE, and re-organise convection around the re-emerging surface low.
%
% The windward-leeward DPD asymmetry (Section~\ref{subsec:ts_dpd_asymmetry})
% provides the boundary-layer moisture perspective on this process:
% BOE-candidate
% systems show smaller leeward DPD increase during terrain crossing, confirming
% that surface evapotranspiration is actively counteracting the
% orographic drying
% on the leeward side.
%
% \section{Surface Energy Partitioning and the Bowen Ratio}
% \label{sec:disc_bowen}
%
% The Bowen ratio trajectory (Section~\ref{subsec:ts_bowen}) is the
% most accessible single-variable summary of the post-landfall
% thermodynamic fate of an MD. The divergence between BOE-candidate and
% non-BOE systems is rapid and quantitatively large: a $\beta$
% difference of $\sim 0.6$-$0.8$ between the two populations after
% 24\,h of land transit represents a fundamentally different surface
% energy regime. The ERA5 HTESSEL land surface model represents this
% regime transition through its soil hydraulic parameterisation, and
% the consistency between the LHF, SHF, \texttt{swvl1}, and DPD time
% series provides confidence that the model is capturing the physically
% relevant transition correctly at the synoptic scale, even if the
% absolute magnitudes are uncertain (Section~\ref{sec:disc_limitations}).
%
% % \section{Implications for Operational Forecasting}
% % \label{sec:disc_forecasting}
% %
% % \yellownote{FIX: refrences}
% % The ASMI-duration relationship (Section~\ref{subsec:asmi_duration}) and the
% % BOE LHF threshold behaviour (Section~\ref{subsec:boe_identified}) together
% % suggest that a simple two-step operational indicator can be defined:
% %
% % \begin{enumerate}
% %   \item \textbf{Pre-landfall ASMI flag}: if the 7-day ASMI over the expected
% %     inland track corridor exceeds $0.35$\,m$^3$\,m$^{-3}$, issue an alert
% %     for possible anomalous inland persistence.
% %   \item \textbf{Post-landfall LHF monitoring}: if the ERA5T (near-real-time)
% %     inner-core LHF remains $\geq 70$\,W\,m$^{-2}$ for $\geq 12$\,h during
% %     the land phase, confirm BOE candidacy and revise the inland rainfall
% %     forecast accordingly.
% % \end{enumerate}
% %
% % This two-step procedure is consistent with the operational framework
% % advocated by \autocite{lodh2022improving}, who demonstrated that soil moisture
% % assimilation improves NWP post-landfall MD prediction, and
% \autocite{osuri2020},
% % who showed that correct soil moisture initialisation is a primary
% determinant
% % of MD simulation skill.
%
% % \section{Scope and Limitations}
% % \label{sec:intro_scope}
% %
% % ERA5 reanalysis (0.25$^\circ$, hourly) is the primary dataset; IMD
% % monsoon reports provide observational validation. The study covers
% % June-September for 2001-2025. At ERA5 spatial resolution ($\approx
% % 28$\,km), mesoscale inner-core convective structure is not
% % resolved, so results characterise synoptic-scale thermodynamic evolution.
%
% \section{Limitations and Sources of Uncertainty}
% \label{sec:disc_limitations}
%
% A systematic assessment of the uncertainties in this study is essential for
% correctly interpreting the magnitude and generality of the results.
%
% \subsection{ERA5 Soil Moisture Overestimation over India}
% \label{subsec:lim_swvl1_bias}
%
% The ERA5 land surface model (HTESSEL) has been shown to
% systematically \emph{overestimate} near-surface soil moisture
% (\texttt{swvl1}) over the Indian subcontinent during the monsoon
% season. \autocite{osuri2020} noted that ERA5 HTESSEL soil moisture over
% the Gangetic plains and Odisha tends to be biased high relative to
% in-situ measurements and SMAP satellite retrievals during
% July-August. The primary causes are:
% \begin{enumerate}
%   \item Overestimation of monsoon precipitation in ERA5
%     \autocite{hersbach2020era5}, which directly wets the soil;
%   \item Underestimation of drainage rates in the HTESSEL hydraulic
%     scheme over clay-loam soils; and
%   \item The absence of explicit irrigation in the model,
%     paradoxically leading to a different type of bias the model does
%     not represent irrigation-enhanced soil moisture but may partially
%     compensate through precipitation overestimation.
% \end{enumerate}
%
% \textbf{Implications for this study.} If \texttt{swvl1} is
% systematically overestimated by $\Delta\theta \approx
% 0.03$-$0.06$\,m$^3$\,m$^{-3}$ during July-August (consistent with the
% bias range reported by \citealt{osuri2020}), then the soil moisture threshold
% of $0.35$\,m$^3$\,m$^{-3}$ identified in this study as a predictor of
% BOE candidacy may be biased high by a similar amount. The \emph{true}
% physical threshold for BOE-enabling soil moisture may therefore be
% closer to $0.29$-$0.32$\,m$^3$\,m$^{-3}$ in terms of actual soil moisture.
% Similarly, the LHF values from ERA5 may be overestimated if the model
% is presenting unrealistically moist soils that evaporate at
% artificially high rates.
% The BOE LHF signal identified here should therefore be treated as an
% upper bound on the actual land-phase LHF, with the true signal
% potentially 10-20\,\% lower.
%
% \textbf{Mitigation.} The use of relative comparisons (BOE candidates
% vs non-BOE systems, ocean phase vs land phase) partially mitigates
% the impact of the absolute bias, as both populations are affected by
% the same systematic overestimation.
% The discrimination between long-lived and short-lived systems, and
% between BOE and non-BOE candidates, is based on differences in LHF
% and \texttt{swvl1} rather than their absolute values, making the
% classification methodology less sensitive to the bias than an
% approach that relies on absolute thresholds.
%
% \subsection{ERA5 Spatial Resolution and Inner-Core Representation}
% \label{subsec:lim_resolution}
%
% At $0.25^\circ$ ($\approx 28$\,km), ERA5 cannot resolve the
% mesoscale convective
% organisation of the MD inner core. The innermost radial bin of the
% Hovm\"{o}ller
% ($r < 25$\,km) contains at most one ERA5 grid point at any time step, and the
% values at this bin are single-point interpolations rather than
% spatial averages.
% Results for $r < 50$\,km should be interpreted as indicative of
% inner-core scale
% behaviour rather than as precise spatial measurements. In particular, the
% precipitation values in the inner-core bin may underestimate peak convective
% rainfall by a factor of 2-3 relative to the actual rainfall intensity at those
% scales, based on comparisons between ERA5 and GPM IMERG in the Indian monsoon
% context \autocite{hersbach2020era5}.
%
% \subsection{ERA5 Accumulated Flux De-Accumulation Errors}
% \label{subsec:lim_deaccum}
%
% The ERA5 LHF and SHF are provided as accumulated values in J\,m$^{-2}$, and
% instantaneous fluxes are derived by differencing consecutive
% accumulation steps
% (Equation~\ref{eq:deaccumulate}). This de-accumulation introduces noise at the
% hourly timescale because the assimilation system generates a new model run at
% each assimilation cycle (every 12\,h), and the accumulated flux is
% reset at the
% start of each run. Values in the first hour of each 12-h assimilation window
% may be affected by initialisation shock and should be treated with caution.
% This issue is most significant for the LHF values near the landfall
% transition,
% which may coincide with an assimilation cycle boundary.
%
% \subsection{SRTM Terrain Representation}
% \label{subsec:lim_srtm}
%
% The SRTM 1-arc-second DEM was acquired in February 2000 (northern winter) and
% represents a point-in-time surface elevation. While the Eastern Ghats bedrock
% elevation has not changed materially since 2000, the SRTM elevation model
% includes the canopy height of forests in the original data, which can add
% 10-20\,m to the ``terrain'' elevation in forested hill areas. For the bulk
% of the terrain analysis (identifying elevation zones of 200-1000\,m), this
% offset is negligible. However, for very precise calculations involving
% slope and roughness, the canopy contribution introduces sub-grid uncertainty
% in the 30-m DEM that is smoothed out when interpolated to track-point
% resolution.
%
% \subsection{Sample Size and Statistical Power}
% \label{subsec:lim_sample}
%
% \yellownote{TODO: update values}
% The 24-year record yields [N$_\text{total}$] confirmed MDs, of which
% [N$_\text{land}$] are land-falling and [N$_\text{BOE}$] qualify as
% BOE candidates. With a BOE-candidate sample of order 10-20 systems,
% the antecedent soil moisture duration correlation and
% Topographic-induced mechanical blocking duration correlation should
% be considered statistically suggestive rather than definitive:
% correlations significant at $p < 0.05$ but not at $p < 0.01$ could
% arise by chance with this sample size. Extension to the full ERA5
% back-extension record (1940-present) or cross-validation with MERRA-2
% would substantially increase statistical power.
%
% \subsection{Land Use, Irrigation, and HTESSEL Representativeness}
% \label{subsec:lim_landuse}
%
% ERA5 HTESSEL does not include explicit representation of agricultural
% irrigation, which is a major contributor to soil moisture over the
% Gangetic plains during the monsoon season \autocite{ankur2022}. In
% reality, irrigation may maintain \texttt{swvl1} at high levels even
% during break periods when natural monsoon rainfall is insufficient.
% The ASMI therefore underrepresents the contribution of irrigation to
% antecedent soil moisture in the most agriculturally active regions.
% Conversely, the absence of irrigation means that ERA5's
% \texttt{swvl1} may underestimate actual soil moisture in the Gangetic
% plains during break periods preceding MD landfall counteracting the
% overestimation bias discussed in Section~\ref{subsec:lim_swvl1_bias}
% and making the net bias direction uncertain in this specific context.
%
% \subsection{Shallow Soil Layer Representation}
% \label{subsec:lim_soil_depth}
%
% This study uses ERA5 \texttt{swvl1} (0-7\,cm), the uppermost soil
% layer, as the primary soil moisture variable. For the short-timescale
% (hourly) LHF diagnostics, \texttt{swvl1} is the most appropriate
% variable as it is most directly coupled to surface evaporation.
% However, for multi-day BOE events lasting $\geq 96$\,h, sustained
% evapotranspiration draws on deeper soil layers (\texttt{swvl2}:
% 7-28\,cm; \texttt{swvl3}: 28-100\,cm). The \texttt{swvl1} layer can
% be dried substantially by the first day of a BOE event's
% precipitation and LHF, even if the deeper layers remain moist enough
% to sustain root-zone evapotranspiration through capillary rise. The
% ASMI using only \texttt{swvl1} may therefore underestimate the total
% moisture reservoir available to sustain BOE in systems lasting $\geq 4$ days.
%
% \subsection{Scope of the Vortex Analysis}
% \label{subsec:lim_vortex}
%
% The three-dimensional dynamical structure of the MD warm/cold core
% evolution, pressure-level vorticity budgets, vertical tilt, IVT, and
% $\theta_e$ vertical profiles was not examined in this thesis. These
% quantities provide a complementary dynamical perspective on the BOE
% and terrain interaction that would strengthen the mechanistic
% attribution. The diagnostic tools for these analyses
% (\texttt{vortex\_analysis.py}) have been implemented and tested as
% part of the codebase but the results are deferred to future work.
%
% \subsection{Fixed Spatial Buffer and Track-Centre Uncertainty}
% \label{subsec:lim_buffer}
%
% The 100\,km spatial buffer used for track-following time series is a
% fixed parameter.
% The optimal buffer scale depends on the size of the MD circulation
% at any given
% time, which varies during the life cycle (contracting as the system decays).
% A fixed buffer may include grid points outside the vortex
% circulation during the
% peak land phase, diluting the inner-core flux signal with background values.
% Sensitivity tests with 50\,km and 200\,km buffers (not shown) reveal that the
% qualitative conclusions are robust, but the quantitative LHF values change by
% $\pm 15$\,\% depending on the buffer choice.
%
% Furthermore, the ERA5 850\,hPa vorticity maximum used to define the
% track centre
% may be displaced 25-75\,km from the actual low-level pressure minimum owing to
% the coarse model resolution and the complex multi-scale structure of the MD
% circulation. This positional uncertainty propagates directly into the
% Hovm\"{o}ller radial coordinate and the time-series buffer extraction.
%
% % \section{Mechanistic Role of Soil Moisture Memory}
% % \label{sec:disc_memory}
% %
% % The persistence of LHF is tied to the soil moisture memory, as
% % evidenced by the coupling efficiency metrics.
% %
% % \yellownote{TODO}
% % \begin{figure}[ht]
% %   \centering
% %   % \includegraphics[width=0.8\textwidth]{output_2024/Global_soil_memory/all_tracks_soil_memory.csv}
% %   % Create a bar or scatter plot from this CSV data
% %   \fbox{\rule{0pt}{7cm}\rule{0.82\textwidth}{0pt}}
% %   \caption{Soil moisture memory scales across the MD ensemble.}
% %   \label{fig:soil_memory}
% % \end{figure}
%
% \section{Historical Context and Prior Work on Indian MD Climatology}
% \label{sec:md_climatology}
%
% \subsection{Comparison of This Study with Prior Methodology}
% \label{subsec:methodology_comparison}
%
% Table~\ref{tab:prior_methods} summarises the methodological
% approaches of prior Indian MD soil moisture studies and positions the
% present study relative to them.
%
% \begin{table}[H]
%   \centering
%   \caption[Comparison of prior Indian MD soil moisture
%   studies]{Summary of prior Indian MD antecedent moisture and
%     post-landfall studies. ``Direct LHF'' indicates whether the study
%   measured surface latent heat flux directly rather than using a proxy.}
%   \label{tab:prior_methods}
%   \begin{tabular}{p{3cm}p{3cm}p{2.8cm}p{3.2cm}}
%     \toprule
%     \textbf{Study}                   & \textbf{Method}
%     & \textbf{Direct LHF?} & \textbf{Key Limitation} \\
%     \midrule
%     \autocite{kishtawal2013enhancement} & Antecedent rainfall proxy
%     & No                   & Proxy, fixed grid       \\
%     \autocite{chang2009}                & WRF NWP sensitivity
%     & No (model)           & Single event            \\
%     \autocite{hunt2017effect}           & NWP SM perturbation
%     & No (model)           & Model, not obs          \\
%     \autocite{osuri2020}                & NWP + ASCAT SM
%     & No (model)           & Model, not obs          \\
%     \autocite{vijay2024boe}             & ERA5 case study
%     & Partial              & Single event, 1999      \\
%     \textbf{This study}              & \textbf{Azimuthal
%     Hovm\"{o}ller + ERA5} & \textbf{Yes}  & ERA5 resolution limits \\
%     \bottomrule
%   \end{tabular}
% \end{table}
%
% The primary novelty of the present study-which
% Table~\ref{tab:prior_methods} makes explicit-is the combination of
% (i) \emph{direct} LHF measurement from ERA5 (not a proxy),
% (ii) a \emph{storm-centred} diagnostic (Azimuthal Hovm\"{o}ller) that
% removes the background climatological gradient, and
% (iii) a \emph{systematic} multi-year (2001-2025) climatology rather
% than single-event analysis.
%
% \subsection{Annual Frequency and Seasonal Distribution}
% \label{subsec:md_frequency}
%
% Historically, the Bay of Bengal generates approximately 5-7 Monsoon
% Depressions per JJAS season, with peak activity in July and August
% when the monsoon trough is at its most pronounced and BoB SSTs are at
% their seasonal maximum of 28-30$^\circ$C \autocite{sikka1977,
% rashid2005track}. The annual count is modulated by large-scale climate modes:
%
% \begin{itemize}
%   \item \textbf{ENSO:} El Ni\~no years typically suppress Indian
%     monsoon rainfall and reduce BoB LPS frequency, as the Walker
%     circulation weakens moisture flux over the BoB. La Ni\~na years
%     enhance LPS genesis.
%   \item \textbf{Indian Ocean Dipole (IOD):} Positive IOD (warm
%     western, cool eastern Indian Ocean) suppresses BoB moisture flux
%     and MD genesis; negative IOD enhances it.
%   \item \textbf{Madden-Julian Oscillation (MJO):} During the active
%     (convective) phase of the MJO over the Indian Ocean, background
%     moisture and vorticity over the BoB are enhanced, increasing MD
%     genesis probability.
% \end{itemize}
%
% \subsection{Landfall Statistics and Preferential Track Corridors}
% \label{subsec:md_landfall_stats}
%
% The majority of BoB MDs make landfall along the northern Andhra
% Pradesh-Odisha-West Bengal coastline, between approximately
% 82$^\circ$E and 88$^\circ$E longitude \autocite{rashid2005track}. This
% clustering reflects the combined effect of the monsoon trough axis
% (which directs systems toward northern Odisha) and the geometry of
% the BoB (which is wider in the south, giving systems forming farther
% south a longer ocean transit that brings them to the northern coast).
% Systems making landfall south of approximately 16$^\circ$N encounter
% the main Eastern Ghats range at higher topographic sensitivity, while
% systems making landfall north of 20$^\circ$N may partially circumvent
% the Ghats by tracking through the Mahanadi river-valley gap.
%
% % \yellownote{TODO: Discussions}
% % The Discussion on:
% % - Comparison of LHF values (70-120 W/m2) with BOE cases
% % (>200 W/m2) and physical explanation (monsoon cloud cover reduces
% % available radiation)
% % - Comparison of ASMI threshold (0.35 m3/m3) with Kishtawal (2013) proxy
% % - Explanation of why the 7-day lag outperforms the 14-day lag (soil
% % moisture memory timescale; clay loam retains moisture 3-7 days)
% % - Implications for operational forecasting
